\documentclass[11pt,letterpaper]{report}

\usepackage[utf8]{inputenc}
\usepackage[T1]{fontenc}
\usepackage[spanish,provide=*]{babel}
\usepackage{amsmath,amssymb}
\usepackage{booktabs}
\usepackage{array}
\usepackage{graphicx}
\usepackage{float}
\usepackage{caption}
\usepackage[margin=2.5cm]{geometry}
\usepackage{enumitem}

\newcommand{\Prob}{\operatorname{Prob.}}
\renewcommand{\arraystretch}{1.15}

\begin{document}

\begin{center}
{\LARGE\bfseries\linespread{1.3}\selectfont 6. ANÁLISIS DE LOS RESULTADOS DE LA SIMULACIÓN\par}
\end{center}
\vspace{0.5cm}

\setcounter{chapter}{6}
\setcounter{page}{107}

%==========================================================
Los resultados de una simulación evalúan la eficiencia u operación de un
sistema. Sin embargo, la confiabilidad de estos resultados dependen del
número de observaciones que se tienen, es decir, la confiabilidad de los
resultados obtenidos en un estudio de simulación, depende de la longitud
de la corrida y del número de corridas. Por consiguiente, es necesario
determinar intervalos de confianza para las variables del sistema que se
están analizando.

La determinación de intervalos de confianza depende de si los resultados
de la simulación son independientes o están correlacionados. Por
ejemplo, en la simulación de un sistema de colas, el tiempo de espera de
un cliente en el sistema, depende del número de clientes que este último
cliente encontró al llegar al sistema (correlación entre $t_i$ y
$t_{i-1}$). También, la obtención de intervalos de confianza depende de
si el sistema se analiza en estado transiente o en estado estable. Por
consiguiente, el objetivo de este capítulo en analizar la determinación
de intervalos de confianza para las dos situaciones mencionadas, es
decir, para: 1) Sistemas cuyos resultados son independientes uno del
otro y 2) Sistemas transientes con alta correlación en sus resultados.
Para el primer caso se utilizan los métodos tradicionales de intervalos
de confianza. Para el segundo caso se utiliza simulación regenerativa.

%==========================================================
\section*{6.1. Métodos de estimación}

Cuando los sistemas que se analizan se encuentran en estado estable, y
además los valores que toman las variables que interactúan en el
sistema son independientes, entonces es posible utilizar los
procedimientos tradicionales de intervalos de confianza. En la
determinación de estos intervalos de confianza se utiliza el teorema
del límite central. Este teorema establece que la suma de $n$ variables
aleatorias independientes, obtenidas de un universo con media $\mu$ y
varianza $\sigma^2$, es aproximadamente una distribución normal con
media $n\mu$ y varianza $n\sigma^2$. También, es obvio que a partir de
este teorema la distribución de:

\begin{equation}
\bar{x} = \frac{1}{n}X_1 + \frac{1}{n}X_2 + \ldots + \frac{1}{n}X_n
\end{equation}

es una distribución normal con media $\mu$ y varianza $\sigma^2/n$ y
por consiguiente:

\begin{equation}
\frac{\bar{x} - \mu}{\sigma/\sqrt{n}}
\end{equation}

sigue una distribución normal estándar con media 0 y varianza 1. De
acuerdo a esta última expresión, y suponiendo un nivel de significado
$\alpha$, el intervalo de confianza del parámetro $\mu$ sería:

\[
\Prob\left\{ -Z_{\alpha/2} \le \frac{\bar{x}-\mu}{\sigma/\sqrt{n}} \le Z_{\alpha/2} \right\} = 1-\alpha
\]

\begin{equation}
\Prob\left\{ \bar{x} - \frac{\sigma}{\sqrt{n}}Z_{\alpha/2} \le \mu \le \bar{x} + \frac{\sigma}{\sqrt{n}}Z_{\alpha/2} \right\} = 1-\alpha
\end{equation}

La expresión anterior establece que el ancho del intervalo de confianza
es:

\begin{equation}
\frac{2\sigma Z_{\alpha/2}}{\sqrt{n}}
\end{equation}

lo cual significa que el tamaño del intervalo sería más pequeño entre
mayor sea el número de observaciones que se tienen de la variable
aleatoria $X$. Por ejemplo, para reducir el intervalo de confianza a la
mitad, es necesario de acuerdo a la expresión anterior, aumentar cuatro
veces el número de observaciones.

En las expresiones anteriores se asume que la varianza es conocida. Sin
embargo, en muchos problemas reales la varianza generalmente se
desconoce. Para este tipo de situaciones, es posible utilizar el
siguiente estimador de la varianza:

\begin{equation}
S^2 = \frac{1}{n-1} \sum_{i=1}^{n} (X_i - \bar{x})^2
\end{equation}

Ahora, si esta identidad se sustituye en la expresión de $Z$, lo que
resulta es:

\begin{equation}
\frac{\bar{x} - \mu}{S/\sqrt{n}}
\end{equation}

lo cual sigue una distribución $t$ con $n-1$ grados de libertad. De
acuerdo a esta nueva distribución y suponiendo un nivel de significado
$\alpha$, el intervalo de confianza del parámetro $\mu$ sería:

\begin{equation}
\Prob\left\{ \bar{x} - \frac{S}{\sqrt{n}}t_{n-1,\alpha/2} \le \mu \le \bar{x} + \frac{S}{\sqrt{n}}t_{n-1,\alpha/2} \right\} = 1 - \alpha
\end{equation}

el cual define un ancho del intervalo de confianza de:

\begin{equation}
\frac{2S\,t_{n-1,\alpha/2}}{\sqrt{n}}
\end{equation}

\subsection*{Ejemplo 6.1}

Suponga que la simulación de un proyecto de inversión arrojó los
resultados que se muestran en la tabla 6.1. Para estos resultados se
desea determinar un intervalo de confianza del 95\% ($\alpha=5\%$) para
la tasa interna de rendimiento (TIR).

Puesto que la varianza del universo de donde procede la muestra se
desconoce, entonces el intervalo de confianza para la TIR, se determina
a partir de la expresión anterior, esto es:

\[
13.37\% \pm \frac{6.71}{\sqrt{100}}(1.99)\%
\]

\[
13.37\% \pm 1.34\%
\]

lo cual define el intervalo $(12.03\%;14.71\%)$ cuyo ancho es de
2.68\%.

\begin{table}[H]
\centering
\caption*{\textbf{Tabla 6.1} Resultados (TIR's) de la simulación de un proyecto de inversión (\%)}
\begin{tabular}{cccccccccc}
\toprule
12 & 8 & 5 & 3 & 7 & 5 & 3 & 29 & 27 & 9 \\
15 & 22 & 7 & 21 & 3 & 15 & 26 & 17 & 8 & 13 \\
18 & 13 & 18 & 17 & 5 & 3 & 4 & 7 & 13 & 19 \\
11 & 16 & 13 & 13 & 9 & 8 & 7 & 13 & 15 & 12 \\
13 & 12 & 14 & 14 & 20 & 22 & 11 & 12 & 18 & 11 \\
17 & 4 & 23 & 8 & 15 & 7 & 18 & 19 & 24 & 15 \\
19 & 24 & 12 & 7 & 18 & 15 & 13 & 8 & 21 & 4 \\
9 & 17 & 11 & 15 & 21 & 18 & 15 & 28 & 13 & 3 \\
14 & 19 & 19 & 18 & 4 & 27 & 21 & 4 & 5 & 8 \\
10 & 21 & 25 & 21 & 7 & 2 & 8 & 9 & 7 & 11 \\
\bottomrule
\end{tabular}
\end{table}

Finalmente, conviene señalar que si el ancho del intervalo se quisiera
reducir a 1.34\%, entonces se necesitaría aumentar a 400 el número de
simulaciones de la TIR.

%==========================================================
\section*{6.2. Simulación regenerativa}

La simulación de un sistema estocástico puede ser visto como un
experimento estadístico, donde primero se construye el modelo del
sistema, el cual captura la esencia del problema y describe su
estructura fundamental en términos de relaciones causa-efecto. Enseguida,
con la ayuda de una computadora se realizan experimentos y se analizan
los resultados de la simulación con el propósito de hacer inferencias
sobre el comportamiento del sistema. Por consiguiente, los experimentos
de simulación son experimentos estadísticos, y como tales, están basados
en procedimientos estadísticos profundos. En particular, las técnicas de
estimación son muy útiles en estudios de simulación, pues permiten hacer
inferencias estadísticas sobre las variables estocásticas analizadas en
el modelo.

Recientemente, una nueva metodología estadística ha sido desarrollada
para analizar los resultados de una simulación. A esta nueva técnica se
le conoce con el nombre de ``Simulación regenerativa''. Este nuevo
enfoque se basa en el hecho de que muchos sistemas estocásticos se
regeneran (empiezan en el mismo punto de partida) a intervalos regulares
de tiempo. Esta particularidad permite al analista observar y estudiar
durante el curso de la simulación, bloques de datos independientes e
idénticamente distribuidos. Más generalmente, la simulación regenerativa
resuelve los problemas de: 1) ¿Cómo empezar la simulación? 2) ¿Cuándo
empezar a colectar datos? 3) Tamaño de la corrida y número de corridas y
4) ¿Cómo manejar los resultados cuando estos están correlacionados?

\subsection*{6.2.1. Simulación de un sistema de colas de un servidor}

La representación de un sistema de colas de un servidor aparece en la
figura 6.1. En esta figura se puede apreciar que los clientes son
generados por una fuente alimentadora, la cual genera $X$ número de
clientes por unidad de tiempo. Si un cliente que llega al sistema
encuentra desocupado al servidor, entonces este cliente empieza a ser
inmediatamente servido y su salida del sistema empieza inmediatamente
después de que su servicio ha terminado. Por el contrario, si un
cliente al arribar encuentra al servidor ocupado, entonces este cliente
pasa a un salón de espera donde esperará su turno para ser servido.
Finalmente, los clientes que permanecen en la cola, son servidos en base
a atender primero a los que llegaron primero a la cola.

\begin{figure}[H]
\centering
\caption*{\textbf{Figura 6.1} Sistema de colas de un servidor}
\fbox{\parbox{0.85\textwidth}{\centering\vspace{1.5cm}
[Figura: Fuente alimentadora $\to$ Clientes $\to$ (Sistema de colas: Salón de espera $\to$ Servidor) $\to$ Clientes servidos]
\vspace{1.5cm}}}
\end{figure}

Para el sistema particular que se está analizando, suponga que la
fuente alimentadora genera un nuevo cliente cada 2 minutos. También,
suponga que el tiempo de servicio es una variable aleatoria
uniformemente distribuida entre 1 y 2 minutos. Además, suponga que la
forma de evaluar el funcionamiento y operación del sistema, es en base
al valor esperado del tiempo que un cliente permanece en la cola
$E(W)$, cuando el sistema se encuentra en una situación de estado
estable. Para este sistema en particular, no existen procedimientos
analíticos capaces de obtener $E(W)$. Por consiguiente, la única forma
de evaluar y analizar este sistema es a través de simulación.

Por consiguiente, el objetivo de la simulación de este sistema, es
analizar los resultados obtenidos con el propósito de estimar $E(W)$.
Más generalmente, se desea que la estimación de $E(W)$ sea confiable, es
decir, se requiere que el valor de $E(W)$ pertenezca a un intervalo con
un cierto nivel de confianza, por ejemplo del 95\% ($\alpha=5\%$).

Una forma razonable de estimar $E(W)$ sería: sea $W_1$ el tiempo de
espera en la cola del cliente 1, $W_2$ el tiempo de espera en la cola
del cliente 2, y así sucesivamente. Por consiguiente, si durante la
simulación del sistema entraron a éste $N$ clientes, entonces el
promedio muestral sería:

\begin{equation}
\frac{W_1 + W_2 + \ldots + W_n}{N}
\end{equation}

el cual representa a un estimador consistente de $E(W)$. Sin embargo, el
promedio muestral sería en general un estimador sesgado de $E(W)$ debido
a las condiciones iniciales. Por ejemplo, si $W_1$ tiene un valor de
cero, entonces los tiempos de espera de los primeros clientes tienden a
ser pequeños.

La forma tradicional de tratar el sesgo del estimador debido a
condiciones iniciales, es correr el modelo de simulación durante algún
tiempo sin colectar ningún dato, hasta que el sistema de colas haya
alcanzado la condición de estado estable. A partir de este instante, se
empieza a colectar la información. Por ejemplo, se puede simular el
sistema para 1000 clientes y sólo considerar los tiempos de espera de
los últimos 500, y entonces usar:

\begin{equation}
\frac{W_{501} + W_{502} + \ldots + W_{1000}}{500}
\end{equation}

como un estimador de $E(W)$. Sin embargo, esta solución no es muy
recomendable por las siguientes razones: 1) Es muy difícil determinar la
longitud del período de estabilización y 2) Una gran cantidad de tiempo
de computadora es desperdiciado.

Otra de las desventajas que se le atribuyen al procedimiento anterior,
es la alta correlación que existe entre $W_i$ y $W_{i+1}$
independientemente de que el sistema se encuentre o no en estado
estable. Por consiguiente, si $W_1, W_2, \ldots, W_n$ no son
independientes, la teoría estadística clásica sería de poca utilidad en
la determinación del intervalo de confianza del tiempo promedio de
espera en la cola.

Se puede observar de los comentarios anteriores, que el sesgo debido a
condiciones iniciales transientes y a la alta correlación de los datos,
ofrecen serios obstáculos en la determinación de $E(W)$. Por
consiguiente, es necesario buscar y proponer un nuevo procedimiento que
elimine y resuelva los problemas antes mencionados. Afortunadamente, sí
existe tal procedimiento al cual se le llama simulación regenerativa.

Para entender la lógica de esta nueva técnica, suponga que al empezar la
simulación se establece $W_1=0$. En seguida, se empieza a correr el
modelo de simulación y valores del tiempo de espera en la cola empiezan
a ser colectados. La obtención de tales datos suponga que son tal y como
se muestran en la tabla 6.2. En esta tabla se puede observar que los
clientes 1, 4, 8, 9, 14, 17 y 20 fueron muy afortunados, puesto que
encontraron desocupado al servidor al momento de arribar al sistema.
También, en esta tabla se puede apreciar que el servidor está
constantemente ocupado desde que llega el cliente 1 hasta que se va el
cliente 3, en seguida se encuentra ocioso hasta que llega el cliente 4,
y así sucesivamente. Si más de 20 clientes son servidos, entonces el
sistema simulado mostrará este mismo comportamiento, es decir, servidor
ocupado, servidor ocioso, servidor ocupado, servidor ocioso, y así
sucesivamente. Por consiguiente, si al tiempo que transcurre entre
empezar a servir y quedar ocioso, se le considera como un ciclo,
entonces para la corta simulación mostrada en la tabla 6.2, existen 6
ciclos, donde el primer ciclo comprende a los clientes 1, 2 y 3, el
segundo ciclo comprende a los clientes 4, 5, 6 y 7 y así sucesivamente.

\begin{table}[H]
\centering
\caption*{\textbf{Tabla 6.2} Tiempo de espera en la cola}
\begin{tabular}{ccc ccc ccc}
\toprule
No. Cliente & $W$ & Ciclo & No. Cliente & $W$ & Ciclo & No. Cliente & $W$ & Ciclo \\
\midrule
1 & 0 & & 8 & 0 & & 14 & 0 & \\
2 & 10 & 1 & 9 & 0 & 3 & 15 & 16 & 5 \\
3 & 5 & & 10 & 18 & & 16 & 12 & \\
4 & 0 & & 11 & 20 & & 17 & 0 & \\
5 & 25 & 2 & 12 & 30 & 4 & 18 & 25 & 6 \\
6 & 15 & & 13 & 40 & & 19 & 30 & \\
7 & 20 & & & & & 20 & 0 & \\
\bottomrule
\end{tabular}
\end{table}

El séptimo ciclo empieza con el arribo del cliente número 20. Por
consiguiente, un nuevo ciclo es iniciado cuando un cliente llega al
sistema y encuentra al servidor ocioso.

Puesto que cada vez que un cliente llega al sistema y encuentra al
servidor ocioso, es equivalente y está gobernado por la misma estructura
probabilística que cuando llega el primer cliente al sistema, entonces
sería razonable agrupar los resultados de la simulación en bloques de
datos, donde el primer bloque contendrá los tiempos de espera en la
cola de los clientes que forman parte del primer ciclo, el segundo
bloque contendrá los tiempos de espera en la cola de los clientes que
forman parte del segundo ciclo, y así sucesivamente. Para la corta
simulación mostrada en la tabla 6.2, los bloques serían:

\[
(W_1,W_2,W_3),\ (W_4,W_5,W_6,W_7),\ (W_8),\ (W_9,W_{10},W_{11},W_{12},W_{13}),\ (W_{14},W_{15},W_{16}),\ (W_{17},W_{18},W_{19})
\]

De este modo, puesto que cada ciclo es iniciado en las mismas
condiciones, los bloques de datos de ciclos sucesivos son
estadísticamente independientes e idénticamente distribuidos. Por
consiguiente, si se define $E_k$ como la suma de los tiempos de espera
de los clientes servidos en el ciclo $k$, y $C_k$ como el número de
clientes servidos en el ciclo $k$, entonces los pares $(E_1,C_1),
(E_2,C_2), (E_3,C_3), (E_4,C_4), (E_5,C_5)$ y $(E_6,C_6)$ son
estadísticamente independientes e idénticamente distribuidos, a pesar de
la alta correlación que existe entre $E_k$ y $C_k$. De la tabla 6.2, los
valores de $(E_k,C_k)$ que se obtienen son los siguientes:

\begin{align*}
(E_1,C_1) &= (15,3) & (E_4,C_4) &= (108,5) \\
(E_2,C_2) &= (60,4) & (E_5,C_5) &= (28,3) \\
(E_3,C_3) &= (0,1) & (E_6,C_6) &= (55,3)
\end{align*}

Por consiguiente, la alta correlación entre $W_k$ y $W_{k+1}$ ha sido
eliminada y los datos han sido agrupados en bloques estadísticamente
independientes e idénticamente distribuidos.

De acuerdo a la agrupación anterior, y suponiendo que $N$ clientes
fueron servidos durante la simulación, a partir de la cual se formaron
$n$ bloques, la estimación del tiempo de espera en la cola se puede
determinar por medio de la siguiente expresión:

\begin{equation}
E(W) = \frac{W_1+W_2+\ldots+W_N}{N} = \frac{E_1+E_2+\ldots+E_n}{C_1+C_2+\ldots+C_n}
\end{equation}

Sin embargo, la segunda parte de la expresión anterior puede ser escrita
de la siguiente forma:

\begin{equation}
E(W) = \frac{(E_1+E_2+\ldots+E_n)/n}{(C_1+C_2+\ldots+C_n)/n}
\end{equation}

Puesto que $(E_1,E_2,\ldots,E_n)$ y $(C_1,C_2,\ldots,C_n)$ son grupos de
variables aleatorias independientes e idénticamente distribuidas,
entonces por el teorema del límite central, la expresión anterior
converge a:

\begin{equation}
\lim_{n \to \infty} E(W) = E(E)/E(C)
\end{equation}

Por consiguiente, el problema de estimar $E(W)$ se traduce a estimar
$E(E)/E(C)$. Y puesto que $E(E)/E(C)$ puede ser estimado a partir de
$(E_1,C_1), (E_2,C_2), \ldots, (E_n,C_n)$, la estadística clásica puede
ser utilizada para hacer inferencias acerca del valor verdadero de
$E(W)$. Además, a través de la estadística clásica se puede determinar
un intervalo de confianza para $E(W)$.

\bigskip
\noindent\emph{Nota sobre esta transcripción:} en el PDF de origen
consultado (páginas 115--116 del libro), faltan las páginas donde Coss
Bu desarrolla la varianza muestral del cociente $\hat{t}=\bar{E}/\bar{C}$
(términos $S_E^2$, $S_C^2$, $S_{EC}$, la expresión general de $S^2$ y
las ecuaciones 6.14 a 6.27) antes de llegar a la ecuación (6.28) del
ancho del intervalo. Esas dos páginas del libro no están en el PDF
subido a \texttt{referencias/}; si se necesitan para el informe, hay
que volver a escanear/subir esas dos páginas.

el cual define un intervalo cuyo ancho es de:

\begin{equation}
\frac{2S\,Z_{\alpha/2}}{\bar{C}\sqrt{n}}
\end{equation}

Finalmente, conviene mencionar que aunque la variable aleatoria
$E(E)/E(C)$ sigue una distribución de probabilidad $t$ puesto que
$\sigma^2$ se desconoce, la expresión del ancho del intervalo sigue
siendo válida ya que para valores de $n$ muy grandes, la distribución
$t$ se aproxima a la distribución normal.

\subsection*{Ejemplo 6.2}

Para los tiempos de espera en la cola mostrados en la tabla 6.2,
determine un intervalo de confianza del 90\% ($\alpha=10\%$) para
$E(W)$.

Para los tiempos de espera mostrados en la tabla 6.2, los bloques de
datos son:

\begin{align*}
(E_1,C_1) &= (15,3) & (E_4,C_4) &= (108,5) \\
(E_2,C_2) &= (60,4) & (E_5,C_5) &= (28,3) \\
(E_3,C_3) &= (0,1) & (E_6,C_6) &= (55,3)
\end{align*}

De acuerdo a estos bloques de datos, los promedios muestrales $\bar{E}$
y $\bar{C}$ y el parámetro $\hat{t}$ serían:

\begin{align*}
\bar{E} &= \frac{15+60+0+108+28+55}{6} = 44.33 \\
\bar{C} &= \frac{3+4+1+5+3+3}{6} = 3.16 \\
\hat{t} &= \frac{\bar{E}}{\bar{C}} = \frac{44.33}{3.16} = 14.03
\end{align*}

Por otra parte, el valor correspondiente de $S_E^2$ sería:

\begin{align*}
S_E^2 &= \frac{1}{5}\Big( (15-44.33)^2 + (60-44.33)^2 + (44.33)^2 \\
&\qquad + (108-44.33)^2 + (28-44.33)^2 + (55-44.33)^2 \Big) \\
S_E^2 &= 1501.06
\end{align*}

(los valores de $S_C^2$, $S_{EC}$ y $S^2$ requeridos para completar el
intervalo de confianza dependen de las ecuaciones de las páginas
faltantes señaladas arriba).

%==========================================================
\section*{PROBLEMAS}

\begin{enumerate}
\item[6.1.] Para el sistema de colas descrito en el ejemplo de este
capítulo, si se simula la llegada de 50 clientes y se obtienen los
tiempos de espera correspondientes, determine un intervalo de confianza
del 95\% para $E(W)$ utilizando el método de simulación regenerativa.

\item[6.2.] Para los datos de interferencia (probabilidad de
interferencia) que se muestran a continuación, determine un intervalo
de confianza del 95\% para la probabilidad de interferencia.

\begin{table}[H]
\centering
\begin{tabular}{cccccccccc}
0.30 & 0.22 & 0.25 & 0.33 & 0.30 & 0.38 & 0.29 & 0.31 & 0.32 & 0.34 \\
0.28 & 0.34 & 0.27 & 0.31 & 0.35 & 0.37 & 0.39 & 0.35 & 0.43 & 0.28 \\
0.29 & 0.33 & 0.23 & 0.36 & 0.25 & 0.35 & 0.30 & 0.27 & 0.28 & 0.24 \\
0.33 & 0.31 & 0.28 & 0.24 & 0.29 & 0.36 & 0.41 & 0.38 & 0.34 & 0.29 \\
0.35 & 0.34 & 0.24 & 0.29 & 0.24 & 0.33 & 0.38 & 0.36 & 0.38 & 0.32 \\
0.37 & 0.29 & 0.37 & 0.22 & 0.32 & 0.28 & 0.42 & 0.41 & 0.37 & 0.33 \\
0.35 & 0.26 & 0.36 & 0.39 & 0.31 & 0.26 & 0.43 & 0.44 & 0.39 & 0.39 \\
0.34 & 0.27 & 0.35 & 0.34 & 0.34 & 0.29 & 0.38 & 0.29 & 0.40 & 0.42 \\
0.28 & 0.31 & 0.31 & 0.33 & 0.35 & 0.34 & 0.39 & 0.30 & 0.42 & 0.31 \\
0.39 & 0.34 & 0.34 & 0.30 & 0.34 & 0.36 & 0.43 & 0.33 & 0.44 & 0.25 \\
\end{tabular}
\end{table}

\item[6.3.] Determine un intervalo de confianza del 95\% para el valor
presente neto del problema 8 del capítulo anterior.

\item[6.4.] Para el problema 9 del capítulo anterior, simule 100 TIR's
y determine un intervalo de confianza del 90\% para la tasa interna de
rendimiento.

\item[6.5.] Para los tiempos de espera en la cola mostrados a
continuación, determine un intervalo de confianza del 90\% para
$E(W)$.

\begin{table}[H]
\centering
\begin{tabular}{ccc ccc ccc}
\toprule
No. Cliente & $W$ & & No. Cliente & $W$ & & No. Cliente & $W$ & \\
\midrule
1 & 0 & & 11 & 0 & & 21 & 15 & \\
2 & 10 & & 12 & 15 & & 22 & 20 & \\
3 & 12 & & 13 & 0 & & 23 & 0 & \\
4 & 8 & & 14 & 0 & & 24 & 10 & \\
5 & 0 & & 15 & 18 & & 25 & 18 & \\
6 & 0 & & 16 & 14 & & 26 & 0 & \\
7 & 18 & & 17 & 0 & & 27 & 0 & \\
8 & 24 & & 18 & 19 & & 28 & 15 & \\
9 & 30 & & 19 & 23 & & 29 & 18 & \\
10 & 28 & & 20 & 25 & & 30 & 0 & \\
\bottomrule
\end{tabular}
\end{table}

\item[6.6.] Si se simula la llegada de 100 clientes a la primera
estación de servicio del sistema de colas descrito en el problema 12
del capítulo anterior, determine un intervalo de confianza del 90\%
para $E(W)$.

\item[6.7.] Si se simula la llegada de 100 clientes al sistema de colas
descrito en el problema 13 del capítulo anterior, determine un
intervalo de confianza del 95\% para $E(W)$.

\item[6.8.] Suponga que en el problema 16 del capítulo anterior, 5
máquinas son asignadas a un mecánico. Si se simula el sistema durante
1000 horas, determine un intervalo de confianza del 90\% para el tiempo
que una máquina permanece ociosa esperando ser reparada.

\item[6.9.] Si se simula la llegada de 100 camiones al sistema de colas
descrito en el problema 17 capítulo anterior, determine un intervalo de
confianza del 99\% para $E(W)$.

\item[6.10.] a) La demanda diaria de un cierto artículo está regida por
una distribución binomial con parámetros $n=9$ y $\theta=1/3$. El
tiempo de entrega es instantáneo. El costo de mantener una unidad en
inventario es de \$25 por día y el costo de faltante es de \$100 por
unidad. Además, el gerente de producción ha establecido como política,
colocar una orden cada vez que el inventario sea menor o igual a 3
unidades. La cantidad ordenada deberá elevar el nivel de inventario a 7
unidades. Para esta política, determine el costo promedio diario en
condiciones de estado estable (Suponga un inventario inicial de 7
unidades y que un ciclo se inicia cada vez que existen 7 unidades al
principio del día).

Sugerencia: El costo promedio diario para este problema se puede
determinar de acuerdo a la siguiente expresión:

\begin{equation}
CT = 25\,E\{(x-D)^+\} + 100\,E\{(D-x)^+\}
\end{equation}

donde:
\begin{align*}
CT &= \text{Costo total diario.} \\
x &= \text{Nivel de inventario al final del día.} \\
D &= \text{Demanda diaria.}
\end{align*}

y puesto que:
\begin{align*}
a^+ &= \max(a,0) \\
(a-b)^+ &= a - \min(a,b)
\end{align*}

entonces, la expresión anterior se transforma en:

\begin{equation}
CT = 25\,E\{x\} + 100\,E\{D\} - 125\,E\{\min(x,D)\}
\end{equation}

b) Si se simulan 30 días y se obtienen los resultados que se muestran a
continuación, determine un intervalo de confianza del 90\% para el
costo promedio diario.
\end{enumerate}

\end{document}
